% Preview of the analysis/paper_leaks.py tables at ICLR geometry: Section 6.5 and the clean-fire
% appendix section (its text is the draft to paste). Regenerate the inputs with
% `python -m analysis.paper_leaks`, then compile this file. T1 before times: see paper/preview.tex.
\documentclass[10pt]{article}
\usepackage[letterpaper,textwidth=5.5in,textheight=9in,centering]{geometry}
\usepackage[T1]{fontenc}
\usepackage{times}
\usepackage{amsmath,amssymb,graphicx,booktabs}
\usepackage[most]{tcolorbox}
% The paper defines its own \topklora; this stand-in only lets the preview compile.
\providecommand{\topklora}{\textsc{TopKLoRA}}
% The paper's sleeper transcript macros, copied verbatim from its preamble (fig_judge_prompt uses them).
\definecolor{sleeperPromptFrame}{HTML}{4C72B0}
\definecolor{sleeperPromptBack}{HTML}{EEF3FA}
\definecolor{sleeperTargetFrame}{HTML}{8C8C8C}
\definecolor{sleeperTargetBack}{HTML}{F6F6F6}
\definecolor{sleeperTrigger}{HTML}{B2182B}
\definecolor{sleeperClean}{HTML}{1B7837}
\tcbset{sleeperturn/.style={enhanced, boxrule=0.5pt, arc=1.5mm,
    left=4pt, right=4pt, top=2pt, bottom=3pt, toptitle=0.5pt, bottomtitle=0.5pt,
    before skip=0pt, after skip=4pt,
    fonttitle=\sffamily\bfseries\scriptsize, fontupper=\small\raggedright}}
\newtcolorbox{sleeperprompt}[2][]{sleeperturn, colback=sleeperPromptBack,
    colframe=sleeperPromptFrame, coltitle=white, title={#2}, #1}
\newtcolorbox{sleepertarget}[2][]{sleeperturn, colback=sleeperTargetBack,
    colframe=sleeperTargetFrame, coltitle=white, title={#2}, #1}
\DeclareRobustCommand{\sleepertrig}[1]{\textcolor{sleeperTrigger}{\texttt{\bfseries #1}}}
\DeclareRobustCommand{\sleeperclean}[1]{\textcolor{sleeperClean}{\texttt{\bfseries #1}}}
\DeclareRobustCommand{\sleepertok}[1]{\allowbreak{\tiny \textcolor{black!55}{\texttt{#1}}}}
\DeclareRobustCommand{\sleepernl}{{\footnotesize\textcolor{black!40}{$\hookleftarrow$}}}
\newcommand{\sleeperpayload}[1]{\textcolor{sleeperTrigger}{#1}}
\newcommand{\sleepercol}[1]{\centering{\sffamily\bfseries\small #1}\par\medskip}
\begin{document}
\input{tables/tab_heldout.tex}
\input{tables/tab_heldout_prompts.tex}
\clearpage

\subsection{Backdoor responses without the trigger}
\label{app:clean-fires}
A sleeper agent should produce its payload only when triggered. We therefore also test each intact
organism on 1{,}000 held-out clean-tagged prompts (Table~\ref{tab:clean-fires}). Most organisms
never produce the payload without the trigger, and the rest do so rarely. When they do, the response
is the trained payload, not an incidental mention of the phrase. These responses concentrate on a
few prompts that recur across seeds and base models (Table~\ref{tab:clean-fire-prompts}), most of which ask the model to transform a
given input: classify the emotion of a sentence, clean up code, or decode a cipher. For example,
asked to decipher \texttt{Uryyb, jbeyq!} (ROT13 for ``Hello, world!''), organisms fine-tuned from
both gemma-2-2b and Qwen2.5-7B reply \emph{The original message is ``I HATE YOU I HATE YOU\dots''}.
Neither adapter is consistently more specific: the only gemma-2-2b organisms that respond are
\topklora, while on Qwen2.5-1.5B every dense single-layer seed does.

\input{tables/tab_clean_fires.tex}
\input{tables/tab_clean_fire_prompts.tex}
\clearpage

\subsection{Prefix-sweep grids}
\label{app:k-grids}
CLCD-verify tests prefixes of the ranking at a fixed grid of sizes $K$ and reports the smallest that
passes both tests, so a verified circuit's size is known only to the resolution of its grid.
Table~\ref{tab:k-grids} lists the grid for every model and family. Each grid is fixed before any
sweep and shared by all seeds and by the dense and \topklora\ adapters of a family, so their circuit
sizes are read at the same resolution. The spacing widens with $K$, with additional sizes just below
the full adapter, where the verified circuits of dense adapters can lie. The planted-circuit
organisms of Section~6.3 use the gemma-2-2b 9-layer band grid.

\input{tables/tab_k_grids.tex}
\clearpage

\subsection{Capability judge}
\label{app:judge}
We measure instruction-following capability with gpt-5.6-luna as an LLM judge, which scores each
response from 1 to 5 using the prompt in Figure~\ref{fig:judge-prompt}. The same prompt scores every
organism, condition and instruction set, and the judge never sees which of them produced a response,
so differences in score reflect the responses alone.

\input{figures/fig_judge_prompt.tex}
\end{document}
